\section{Conclusion}
We presented Kinema4D, a novel framework that shifts the paradigm of robotic simulation to 4D spatial-temporal reasoning. By integrating a kinematics-driven grounding stage with a diffusion transformer-based generative pipeline, we effectively decouple deterministic robot motion from stochastic environmental reactions. Experimental results demonstrate that Kinema4D can generalize to simulate diverse real-world dynamics. By providing a 4D world, Kinema4D paves the way for scalable, high-fidelity, and complex embodied simulations.

\paragraph{Limitations.} Since the environmental dynamics are learned through statistical synthesis rather than governed by explicit physical constraints (such as rigid-body dynamics or friction coefficients), our method may occasionally produce behaviors that violate conservation laws or exhibit penetration artifacts, leaving a future research direction of incorporating physical laws into our model.
% While our 4D-based control and modeling significantly mitigates visual “hallucination”, the lack of a symbolic physics solver means the framework prioritizes visual and geometric plausibility over strict analytical correctness, leaving a future research direction of incorporating physical laws into our model.